% =====================================================================
%  M3 -- Two variables, contours and the zig-zag
% =====================================================================
\subsection{M3: Two variables and the zig-zag}
\setprob{M3}

\begin{frame}{\probtitle{M3}{Two variables, contours and the zig-zag}}
  \begin{probbox}[Problem M3 \hfill Mathematical \quad $\star\star$]
    Let $J(x_1, x_2) = x_1^2 + 4x_2^2$ and start from $\x_0 = (2, 1)$.
    \begin{enumerate}[(a)]
      \item Find $\nabla J$ and take two GD steps with $\alpha = 0.1$. Does $J$ go down?
      \item What is the largest $\alpha$ for which GD still converges? Take two steps with $\alpha = 0.2$.
      \item Repeat (a) on the round bowl $x_1^2 + x_2^2$. Compare the two paths.
    \end{enumerate}
  \end{probbox}
  \footnotesize\textcolor{mGrey}{Practice: the update with several parameters, and how the shape of the contours changes GD.}
\end{frame}

\begin{frame}{\soltitle{M3}{1}{3}}
  \textbf{(a) Two steps with $\alpha = 0.1$.} $\nabla J = (2x_1,\ 8x_2)$, so each coordinate has its own factor:
  \[
    x_1 \leftarrow (1 - 2\alpha)\,x_1 = 0.8\,x_1, \qquad x_2 \leftarrow (1 - 8\alpha)\,x_2 = 0.2\,x_2 .
  \]
  \begin{center}\small
  \begin{tabular}{@{}c c c c@{}}
    \toprule
    step & $\x_k$ & $\nabla J(\x_k)$ & $J(\x_k)$ \\
    \midrule
    0 & $(2,\ 1)$ & $(4,\ 8)$ & $8$ \\
    1 & $(1.6,\ 0.2)$ & $(3.2,\ 1.6)$ & $2.72$ \\
    2 & $\ans{(1.28,\ 0.04)}$ & $(2.56,\ 0.32)$ & $1.6448$ \\
    \bottomrule
  \end{tabular}
  \end{center}
  $J$ falls from $8$ to $1.64$. The steep $x_2$ direction is almost done after two steps, but $x_1$ has only
  lost $36\%$: the flat direction is the slow one.
\end{frame}

\begin{frame}{\soltitle{M3}{2}{3}}
  \begin{columns}[T]
    \begin{column}{0.52\textwidth}
      \small
      \textbf{(b) Largest stable $\alpha$.} Both factors must satisfy $|1 - 2\alpha| < 1$ and $|1 - 8\alpha| < 1$.
      The steep direction is stricter:
      \[
        0 < \alpha < \tfrac{2}{8} = \ans{0.25}
      \]
      With $\alpha = 0.2$ the factors are $0.6$ and $-0.6$:
      \[
        (2, 1)\to(1.2,\ -0.6)\to(0.72,\ 0.36)
      \]
      $x_2$ flips sign every step. That is the \alert{zig-zag} across the valley. It still converges, since $0.6 < 1$.
    \end{column}
    \begin{column}{0.46\textwidth}
      \centering
      \includegraphics[width=\linewidth, height=0.6\textheight, keepaspectratio]{fig_zigzag_scaling}
    \end{column}
  \end{columns}
\end{frame}

\begin{frame}{\soltitle{M3}{3}{3}}
  \textbf{(c) The round bowl $x_1^2 + x_2^2$.} Now $\nabla J = (2x_1, 2x_2)$, and with $\alpha = 0.1$ both factors are $0.8$:
  \[
    (2, 1) \to (1.6,\ 0.8) \to \ans{(1.28,\ 0.64)},\qquad J:\ 5 \to 3.2 \to 2.048 .
  \]
  Every point is a multiple of $(2,1)$: the path is a \alert{straight line} into the minimum, because $-\nabla J$ points
  right at it.
  \begin{keybox}[Take-away]
    On the round bowl any $\alpha < 1$ works and $\alpha = 0.5$ finishes in one step.
    On the stretched bowl the steep direction caps $\alpha$ at $0.25$, and the flat direction converges slowly.
    Equal curvatures speed up GD.
  \end{keybox}
\end{frame}
